% ============================================================
%  THE ORTYX PROTOCOL
%  Appendix G: The Formal Grammar of the Ortyx Protocol
% ============================================================

\chapter{The Formal Grammar of the \Ortyx{} Protocol}
\label{app:grammar}

\noindent
This appendix assembles the complete formal specification
of the \Ortyx{} Protocol in compact notation. It provides
a reference for the audit operators, failure geometries,
classification taxonomy, and the recursive self-application
equation. The prose interpretations appear in Parts~III
and~V; this appendix gives the formal definitions alone.

\section{The System Triple}
\label{app-g:triple}

A theoretical framework under audit is represented as:
\[
  \SysTriple = (\mathcal{A}, \mathcal{M}, \mathcal{E}),
\]
where $\mathcal{A}$ is the axiom set, $\mathcal{M}$
is the mechanism set, and $\mathcal{E}$ is the empirical
claim set. Let $\mathcal{E}_{\mathcal{A}}$ denote the
claims in $\mathcal{E}$ entailed by $\mathcal{A}$ alone.

\section{The Four Audit Operators}
\label{app-g:operators}

\subsection{Structural Consistency}
\[
  \AuditS(\SysTriple)
  = \frac{|\mathcal{E}_{\mathcal{M}} \setminus \mathcal{E}_{\mathcal{A}}|}{|\mathcal{E}|},
\]
the fraction of empirical claims non-trivially derived
from the mechanisms beyond the axioms alone. Range: $[0,1]$.

\subsection{Falsifiability}
\[
  \AuditF(\SysTriple)
  = \frac{|\mathcal{E}_{\mathrm{falsifiable}}|}{|\mathcal{E}|},
\]
where $\mathcal{E}_{\mathrm{falsifiable}}$ is the set
of claims for which a specific disconfirming observation
can be specified. Range: $[0,1]$.

\subsection{Projection Dependence}
\[
  \AuditP(\SysTriple)
  = \frac{1}{|\mathcal{E}|} \sum_{e \in \mathcal{E}}
  \mathbf{1}\!\bigl[\sup_{x \in X_{\mathrm{adm}}}
  \Delta_{\pi_e}(x) > \varepsilon_e\bigr],
\]
where $\Delta_{\pi_e}(x) = |\ConstrX(x) - \ConstrM(\pi_e(x))|$
and $\varepsilon_e$ is the acceptable projection defect
for claim $e$. Range: $[0,1]$.

\subsection{Recursive Closure}
\[
  \AuditR(\SysTriple)
  = \mathbb{E}_T\!\left[\SemClose(t)\right]
  = \mathbb{E}_T\!\left[\frac{\|F(M_t)\|}{\|G(X_t)\|}\right],
\]
where $T$ is the evaluation period. Range: $[0, \infty)$,
normalized to $[0,1]$ for use in $\EpStab$ by an appropriate
monotone transform.

\section{The Epistemic Stability Functional}
\label{app-g:stability}

\[
  \EpStab(\SysTriple)
  = \alpha\,\AuditS(\SysTriple)
  + \beta\,\AuditF(\SysTriple)
  - \gamma\,\AuditP(\SysTriple)
  - \delta\,\AuditR(\SysTriple),
\]
where $\alpha, \beta, \gamma, \delta > 0$ are
context-sensitive weights with standard values
$\alpha = \beta = 0.35$, $\gamma = \delta = 0.15$
for mature frameworks; and $\alpha = 0.2$,
$\beta = 0.5$, $\gamma = 0.2$, $\delta = 0.1$
for early-stage research programmes.

\section{The Four Failure Geometries}
\label{app-g:failures}

\begin{align*}
  \FailGeo{1} &: \text{Decorative Formalism}\\
  &\quad \text{Notation without admissible constraint.}\\[4pt]
  \FailGeo{2} &: \text{Definitional Closure}\\
  &\quad \mathcal{E} \subseteq \mathcal{E}_{\mathcal{A}}
  \text{ (conclusions embedded in definitions).}\\[4pt]
  \FailGeo{3} &: \text{Semantic Universality Collapse}\\
  &\quad \forall O: O \text{ consistent with }
  \SysTriple \text{ (no disconfirming observation).}\\[4pt]
  \FailGeo{4} &: \text{Operational Severance}\\
  &\quad \exists e \in \mathcal{E}: \text{no measurement}
  \text{ procedure for quantities named in } e.
\end{align*}

\section{The Four Inferential Levels}
\label{app-g:levels}

\begin{align*}
  \LevelI  &: \text{Engineering utility.}
  \text{ Measurable performance claims.}\\
  \LevelII &: \text{Computational metaphor.}
  \text{ Structural analogies generating research directions.}\\
  \LevelIII&: \text{Cognitive phenomenology.}
  \text{ Reports on the character of experience.}\\
  \LevelIV &: \text{Ontological claim.}
  \text{ Assertions about what phenomena fundamentally are.}
\end{align*}

\section{The Decomposition and Reinterpretation}
\label{app-g:decomposition}

\textit{Phase 3 (Decomposition):}
\[
  \SysTriple \xrightarrow{\AuditOp} (\ThetaS, \ThetaU),
\]
where $\ThetaS$ contains claims with acceptable
audit scores and $\ThetaU$ contains claims that fail
the audit at significant severity levels.

\textit{Phase 4 (Reinterpretation):}
\[
  \RSVPfunctor : \ThetaS \to \ThetaStar,
\]
a structure-preserving map from survivable components
to a lower-projection-defect expressive context.

\textit{Phase 5 (Self-application):}
\[
  \OrtOp_{n+1} = \AuditOp(\OrtOp_n).
\]
The protocol survives through continuous self-application
rather than through the achievement of a stable fixed point.

\section{Five Reconstruction Categories}
\label{app-g:categories}

After decomposition and prior to reinterpretation,
the Interlude classifies surviving and non-surviving
components into five categories:

\begin{enumerate}
  \item \textit{Discarded}: Not recoverable within any
  reasonable reinterpretation; $\FailGeo{2}$ or $\FailGeo{3}$
  at high severity.
  \item \textit{Unstable but Recoverable}: In $\ThetaU$
  due to missing specification; could migrate to $\ThetaS$
  if gaps are closed.
  \item \textit{Metaphorically Useful}: Survives as
  \LevelII{} analogy; does not survive as \LevelIV{} claim.
  \item \textit{Operationally Grounded}: In $\ThetaS$
  at \LevelI{}; survives the full audit.
  \item \textit{Reconstructible}: In $\ThetaS$;
  translatable via $\RSVPfunctor$ to $\ThetaStar$
  with explicit constraint-faithfulness acknowledgment.
\end{enumerate}

\section{The Constraint Horizon Principle}
\label{app-g:constraint-horizon}

From Chapter~24: designing a system as though computational
resources are severely constrained --- even when they
are not --- produces architectures with lower projection
defect, higher structural consistency, and stronger
reuse properties than designing without explicit constraint.
The formal statement: for any design problem $P$ with
available resources $R$, introduce an artificial
constraint $R' \subset R$ as a design horizon. The
resulting design $D(R')$ will exhibit:
\[
  \AuditP(D(R')) \le \AuditP(D(R)),
\]
\[
  \AuditS(D(R')) \ge \AuditS(D(R)),
\]
for the same specification $\SysTriple$, because constraint
forces explicit specification of admissibility conditions
that unconstrained design can leave implicit.
